% Part: incompleteness
Gödel 不完全性定理断言，对于某些足够强的形式系统，存在该语言中的一些语句，在系统内既不能被证明也不能被证否。这适用于包含了足够算术的系统，例如 Robinson 算术或 Peano 算术。不可判定语句并非只有单一的一个，在这些理论中存在着无穷多个这样的不可判定语句。粗略地说，该定理陈述如下：如果一个形式系统包含了足够算术，是一致的，并且是可公理化的，那么它是不完全的。
% Chapter: representability-in-q
% Section: introduction

\documentclass[../../../include/open-logic-section]{subfiles}

\begin{document}

\olfileid{inc}{req}{int}
\olsection{引言}

不完全性定理适用于那些能够表达并证明有关可计算函数基本事实的理论。
我们将描述一种满足此条件的极小理论，称为“$\Th{Q}$”（有时也称作“Robinson 的 $Q$”，以 Raphael Robinson 命名）。
我们将说明一个函数在 $\Th{Q}$ 中\emph{可表示}是什么意思，然后证明以下定理：

\begin{quote}
  一个函数在 $\Th{Q}$ 中可表示，当且仅当它是可计算的。
\end{quote}

一方面，这为我们提供了可计算性的另一种模型。
另一方面，我们也将利用它，通过将停机问题归约到该集合，来证明集合 $\Setabs{!A}{\Th{Q} \Proves !A}$ 是不可判定的。
最终，我们将证明比这强得多的结论。

$\Th{Q}$ 的语言是算术语言；$\Th{Q}$ 包含以下公理（这些公理需与带等词的一阶逻辑的其余公理和规则配合使用）：
\begin{align*}
& \lforall[x][\lforall[y][(\eq[x'][y'] \lif \eq[x][y])]] \tag{$!Q_1$}\\
& \lforall[x][\eq/[\Obj 0][x']] \tag{$!Q_2$}\\
& \lforall[x][(\eq[x][\Obj 0] \lor \lexists[y][\eq[x][y']])] \tag{$!Q_3$}\\
& \lforall[x][\eq[(x + \Obj 0)][x]] \tag{$!Q_4$}\\
& \lforall[x][\lforall[y][\eq[(x + y')][(x + y)']]] \tag{$!Q_5$}\\
& \lforall[x][\eq[(x \times \Obj 0)][\Obj 0]] \tag{$!Q_6$}\\
& \lforall[x][\lforall[y][\eq[(x \times y')][((x \times y) + x)]]] \tag{$!Q_7$}\\
& \lforall[x][\lforall[y][(x < y \liff \lexists[z][\eq[(z' + x)][y]])]] \tag{$!Q_8$}
\end{align*}
对于每个自然数 $n$，将数字符 $\num{n}$ 定义为项 $\Obj{0}^{\prime\prime\ldots\prime}$，其中总共含有 $n$ 个撇号。
因此，$\num{0}$ 就是常项符号~$\Obj{0}$ 本身，$\num{1}$ 是 $\Obj{0}'$，$\num{2}$ 是 $\Obj{0}''$，依此类推。

作为一种算术理论，$\Th{Q}$ 是\emph{极其}弱的；例如，甚至无法证明像 $\lforall[x][\eq/[x][x']]$ 或 $\lforall[x][\lforall[y][(x + y) = (y +
    x)]]$ 这样非常简单的事实。
但我们会看到，$\Th{Q}$ 之所以如此有趣，很大程度上\emph{正是因为}它如此之弱。事实上，它的强度刚好足以让不完全性定理成立。
$\Th{Q}$ 有趣的另一个原因是它具有\emph{有限}的公理集。

比 $\Th{Q}$ 更强的理论（称为 \emph{Peano 算术} $\Th{PA}$）是通过向 $\Th{Q}$ 添加归纳模式得到的：
\[
(!A(\Obj 0) \land \lforall[x][(!A(x) \lif !A(x'))]) \lif \lforall[x][!A(x)]
\]
其中 $!A(x)$ 是任意公式。如果 $!A(x)$ 包含 $x$ 以外的自由变元，我们就在公式前端添加全称量词以约束所有这些变元（从而使归纳模式的对应实例成为一个语句）。
例如，如果 $!A(x, y)$ 还包含自由变元 $y$，则对应的实例为
\[
\lforall[y][((!A(\Obj 0) \land \lforall[x][(!A(x) \lif !A(x'))]) \lif
  \lforall[x][!A(x)])]
\]
利用归纳模式的实例，从 $\Th{PA}$ 的公理能比从 $\Th{Q}$ 的公理证明多得多的结论。
事实上，要找到关于自然数的、在 Peano 算术中\emph{不能}被证明的“自然”语句，是需要花相当大功夫的！

\begin{defn}
\ollabel{defn:representable-fn}
  一个从自然数到自然数的函数 $f(x_0,\ldots,x_k)$
  称为{\em 在 $\Th{Q}$ 中可表示的}，如果存在公式 $!A_f(x_0,\dots,x_k,y)$
  使得每当 $f(n_0,\dots,n_k) = m$ 时，$\Th{Q}$ 可证明
\begin{enumerate}
\item\ollabel{defn:rep:a} $!A_f(\num{n_0}, \dots, \num{n_k}, \num{m})$
\item\ollabel{defn:rep:b} $\lforall[y][(!A_f(\num{n_0}, \dots,
\num{n_k}, y) \lif \num{m} = y)]$。
\end{enumerate}
\end{defn}

该定义还有其他等价表述；例如，我们可以等价地要求 $\Th{Q}$ 证明 $\lforall[y][(!A_f(\num{n_0}, \dots,
\num{n_k}, y) \liff \eq[y][\num{m}])]$。

\begin{thm}
\ollabel{thm:representable-iff-comp}
一个函数在 $\Th{Q}$ 中可表示，当且仅当它是可计算的。
\end{thm}

证明该定理分为两个方向。一旦语法的算术化建立起来，从左到右的方向是相当直接的。另一个方向则需要更多工作。
基本想法如下：我们选取“广义递归”作为“可计算”的精确化表述，并证明每个广义递归函数都在 $\Th{Q}$ 中可表示。
回顾一下，如果一个函数可以由零函数 $\Zero$、后继函数 $\Succ$ 和投影函数 $\Proj{n}{i}$ 出发，
通过复合、原始递归和正则极小化定义，那么它就是广义递归的。
因此，证明每个广义递归函数都在 $\Th{Q}$ 中可表示的一种方法是：证明基本函数是可表示的，并且只要某些函数是可表示的，
由它们通过复合、原始递归和正则极小化定义出的函数也是可表示的。
换句话说，我们可以证明基本函数是可表示的，并且可表示函数在复合、原始递归和正则极小化下是“封闭的”。
这就保证了每个广义递归函数都是可表示的。

事实证明，证明可表示函数在原始递归下封闭的那一步是很困难的。
为了避免这一步，我们首先表明，实际上可以不使用原始递归。
也就是说，我们证明每个广义递归函数都可以仅从基本函数出发，通过复合和正则极小化来定义。
为此，我们证明原始递归实际上可以通过一个特定的正则极小化来实现。
然而，为了让此方法奏效，我们必须添加一些额外的基本函数：加法、乘法，以及恒等关系~$\Char{=}$ 的特征函数。
然后，我们通过证明\emph{这些}基本函数在~$\Th{Q}$ 中都是可表示的，并且可表示函数在复合和正则极小化下封闭，就能证明该定理。
\end{document}